\section{相关技术基础}

\subsection{大语言模型推理与端云协同}

大语言模型生成通常采用自回归解码方式：在给定系统提示、用户输入和已生成历史前缀的条件下，模型预测下一个词元的概率分布，并按照采样或贪心策略选出当前词元，再将其加入上下文继续下一步生成\cite{vaswani2017attention,touvron2023llama2}。这一过程具有顺序性，因此推理效率同时受单次前向计算速度和每一步解码组织方式影响。

从概率建模角度看，给定输入上下文 $x$ 和输出序列 $y=(y_1,y_2,\ldots,y_T)$，自回归语言模型将条件生成概率分解为
\begin{equation}
  p_{\theta}(y \mid x)
  = \prod_{t=1}^{T} p_{\theta}(y_t \mid x, y_{<t}),
  \label{eq:autoregressive-factorization}
\end{equation}
其中，$p_{\theta}$ 表示目标模型给出的词元分布，$y_{<t}$ 表示第 $t$ 个词元之前已经生成的前缀。式 \eqref{eq:autoregressive-factorization} 中每一步生成都依赖前一步结果，因此大语言模型推理天然具有串行特征。

在工程实现中，为避免对相同前缀重复计算，推理系统通常使用 KV Cache 或等价状态缓存机制存储中间结果。对于长上下文、多轮生成和持续会话场景，缓存管理会直接影响推理效率\cite{kwon2023pagedattention,dao2022flashattention}。因此，本地推理、远程推理和端云协同推理都需要维护运行时状态。

若将第 $t$ 步之前的缓存记为 $C_{t-1}$，则一次增量解码可抽象为
\begin{equation}
  (p_{\theta}(\cdot \mid x,y_{<t}), C_t)
  = F_{\theta}(y_{t-1}, C_{t-1}),
  \label{eq:kv-cache-update}
\end{equation}
其中，$F_{\theta}$ 表示模型前向计算与缓存更新过程。缓存 $C_t$ 的连续性决定了系统能否避免重复计算长前缀。在跨设备投机解码中，草稿侧和目标侧虽然不共享同一份缓存，但必须共享一致的已接受前缀，否则式 \eqref{eq:kv-cache-update} 所隐含的条件上下文将不一致。

端云协同推理的基本思想是在不同计算节点之间分配推理任务，使终端承担轻量任务，远端承担计算量较大的任务\cite{satyanarayanan2017edge,shi2016edge}。在本文研究场景中，Android 移动端承担草稿提案，桌面端承担目标验证，移动端因此参与到生成过程内部。

端云协同下的一轮生成耗时可表示为
\begin{equation}
  T_{\mathrm{step}}
  = T_{\mathrm{draft}} + T_{\mathrm{comm}} + T_{\mathrm{verify}} + T_{\mathrm{sync}},
  \label{eq:collaborative-latency}
\end{equation}
其中，$T_{\mathrm{draft}}$ 为草稿生成耗时，$T_{\mathrm{comm}}$ 为跨设备通信耗时，$T_{\mathrm{verify}}$ 为目标侧验证耗时，$T_{\mathrm{sync}}$ 为状态同步与回滚相关耗时。该分解为后续实验中分析草稿侧开销、验证侧开销和跨设备协同开销提供了统一口径。

\subsection{投机解码基本原理}

投机解码通过引入一个草稿模型和一个目标模型协同完成生成。草稿模型依据当前上下文先提出若干候选词元，目标模型再对这些候选进行检查\cite{leviathan2023speculative,chen2023speculative}。当前若干词元与目标模型在相同前缀下的预期一致时，这些词元被接受；若在某一位置不一致，则系统停止接受，并由目标模型返回修正词元作为修正输出。

设草稿模型分布为 $q(\cdot \mid s)$，目标模型分布为 $p(\cdot \mid s)$，其中 $s$ 表示当前已接受前缀。草稿侧一次提出长度为 $K$ 的候选序列
\begin{equation}
  \hat{y}_{1:K}=(\hat{y}_1,\hat{y}_2,\ldots,\hat{y}_K),
  \quad
  \hat{y}_i \sim q(\cdot \mid s,\hat{y}_{<i}).
  \label{eq:draft-proposal}
\end{equation}
在标准投机采样中，第 $i$ 个候选词元的接受概率可写为
\begin{equation}
  \alpha_i
  =
  \min\left(
    1,
    \frac{p(\hat{y}_i \mid s,\hat{y}_{<i})}
         {q(\hat{y}_i \mid s,\hat{y}_{<i})}
  \right).
  \label{eq:speculative-acceptance}
\end{equation}
由式 \eqref{eq:speculative-acceptance} 可知，草稿模型与目标模型在相同前缀下的条件分布越接近，候选词元被接受的概率越高。因此，词元空间一致性、采样参数一致性和状态同步都会影响投机解码的接受率。

投机解码的核心状态是已接受前缀，即当前已经被目标侧确认的权威前缀。所有后续草稿提案和目标验证都必须围绕这一前缀展开。如果该前缀在草稿侧和目标侧之间不能保持一致，则投机会话将无法保持有效。

若一轮提案中实际接受的草稿词元数量为 $A$，则 $A$ 可表示为
\begin{equation}
  A = \sum_{i=1}^{K}
      \mathbf{1}\{\hat{y}_1,\ldots,\hat{y}_i \text{ 均被接受}\},
  \label{eq:accepted-count}
\end{equation}
其中 $\mathbf{1}\{\cdot\}$ 为指示函数。在每个位置接受事件近似独立且接受概率近似为 $\bar{\alpha}$ 的简化假设下，单轮期望接受长度可近似为
\begin{equation}
  \mathbb{E}[A]
  \approx
  \sum_{i=1}^{K} \bar{\alpha}^{i}.
  \label{eq:expected-accepted-length}
\end{equation}
该式是用于解释实验现象的近似模型，反映两个关系：草稿长度 $K$ 增大不必然带来线性收益；接受率下降会降低单轮有效推进词元数。

从候选组织形式看，投机解码可采用线性提案或树形提案。线性提案一次提出一段词元切片，实现复杂度较低；树形提案在每个深度上保留多个候选分支，为目标验证提供局部结构信息\cite{miao2023specinfer,li2024eagle}。树形提案同时对词元对齐、分支概率和验证逻辑提出要求。

结合式 \eqref{eq:collaborative-latency} 和式 \eqref{eq:expected-accepted-length}，跨设备投机解码的平均吞吐率可近似表示为
\begin{equation}
  R
  \approx
  \frac{\mathbb{E}[A] + 1}{T_{\mathrm{step}}},
  \label{eq:throughput-model}
\end{equation}
其中额外的 $1$ 表示发生拒绝或完成验证后由目标侧给出的修正或后续词元。式 \eqref{eq:throughput-model} 将吞吐率与草稿提案质量、每轮跨设备协同开销联系起来。在端云协同场景中，接受率、通信、状态同步和验证侧运行时开销共同决定加速效果。

\subsection{典型方法与本文切入点}

已有研究将投机解码作为一类系统化的推理加速机制进行讨论\cite{leviathan2023speculative,chen2023speculative,miao2023specinfer,li2024eagle}。Speculative Decoding 和 Speculative Sampling 给出了草稿模型与目标模型协作的基本框架\cite{leviathan2023speculative,chen2023speculative}，SpecInfer 讨论了词元树验证问题\cite{miao2023specinfer}，Medusa 与 REST 分别从多头候选和检索候选角度扩展了草稿来源\cite{cai2024medusa,he2023rest}，EAGLE 强调统一词元空间和输出分布保持问题\cite{li2024eagle}。本文不以提出新的投机解码理论为目标，而是将相关方法应用于 Android 移动端与桌面端协同场景，重点解决以下问题：

\begin{enumerate}[label=\arabic*.]
  \item 如何设计适合跨设备的投机会话生命周期。
  \item 如何在统一协议边界下组织草稿提案、目标验证和状态回退。
  \item 如何在跨设备场景中保持草稿侧与目标侧的真实词元空间一致。
  \item 如何在实际系统环境下分析协同开销与性能瓶颈。
\end{enumerate}
